% GENERATED FILE. Edit canonical lessons, then run npm run generate:books.
% canonical-source-hash: fnv1a64:3c27fe33735d8a5c
% canonical-lessons: JA-C112-nouka, JA-C112-daiku, JA-C112-kashu, JA-C112-shigoto, JA-C112-yasumi

\chapter{Farmhouse, Carpenter, Singer, Work, Holiday}
\label{ch:ja-farmhouse-carpenter-singer-work-holiday}
\hlchaptermodality{112}

\begin{chapteropening}
\textbf{By the end of this chapter:} \emph{I can say a farmhouse, a carpenter, a singer, work and a holiday.}

Complete the last lesson of chapter 112: I can say a farmhouse, a carpenter, a singer, work and a holiday.
\end{chapteropening}

\section[\texorpdfstring{\ja{のうか} (nouka)}{のうか (nouka)}]{\ja{のうか} (nouka) — a farmhouse}

\label{lesson:JA-C112-nouka}



\begin{quote}
Before the new one: say the Japanese for a road, a highway, then the Japanese for a pavement.
\end{quote}

\subsection*{You'll want to know: \ja{のうか}}
\textbf{\ja{のうか}} — \emph{nouka} — \textquotedblleft{}a farmhouse\textquotedblright{}.

\begin{grammarlens}[title={What you've built}]
One more word for this chapter.
\end{grammarlens}

\subsection*{Your turn}
\begin{itemize}
\raggedright
  \item \emph{Say it:} \emph{nouka}
  \item \emph{Say it:} \emph{nouka}, once more
  \item \emph{Say it:} say \emph{nouka}
  \item \emph{Recall:} say the Japanese for a road, a highway, then the Japanese for a pavement, then say \emph{nouka} again
\end{itemize}

\begin{tcolorbox}[breakable,colback=teal!4,colframe=teal!35!black,title={Before you move on}]
What does \ja{のうか} mean? (\textquotedblleft{}A farmhouse\textquotedblright{}.) Say it once more.
\end{tcolorbox}
\section[\texorpdfstring{\ja{だいく} (daiku)}{だいく (daiku)}]{\ja{だいく} (daiku) — a carpenter}

\label{lesson:JA-C112-daiku}



\begin{quote}
Before the new one: say the Japanese for a pavement, then the Japanese for a farmhouse.
\end{quote}

\subsection*{You'll want to know: \ja{だいく}}
\textbf{\ja{だいく}} — \emph{daiku} — \textquotedblleft{}a carpenter\textquotedblright{}.

\begin{grammarlens}[title={What you've built}]
One more word for this chapter.
\end{grammarlens}

\subsection*{Your turn}
\begin{itemize}
\raggedright
  \item \emph{Say it:} \emph{daiku}
  \item \emph{Say it:} \emph{daiku}, once more
  \item \emph{Say it:} say \emph{daiku}
  \item \emph{Recall:} say the Japanese for a pavement, then the Japanese for a farmhouse, then say \emph{daiku} again
\end{itemize}

\begin{tcolorbox}[breakable,colback=teal!4,colframe=teal!35!black,title={Before you move on}]
What does \ja{だいく} mean? (\textquotedblleft{}A carpenter\textquotedblright{}.) Say it once more.
\end{tcolorbox}
\section[\texorpdfstring{\ja{かしゅ} (kashu)}{かしゅ (kashu)}]{\ja{かしゅ} (kashu) — a singer}

\label{lesson:JA-C112-kashu}



\begin{quote}
Before the new one: say the Japanese for a farmhouse, then the Japanese for a carpenter.
\end{quote}

\subsection*{You'll want to know: \ja{かしゅ}}
\textbf{\ja{かしゅ}} — \emph{kashu} — \textquotedblleft{}a singer\textquotedblright{}.

\begin{grammarlens}[title={What you've built}]
One more word for this chapter.
\end{grammarlens}

\subsection*{Your turn}
\begin{itemize}
\raggedright
  \item \emph{Say it:} \emph{kashu}
  \item \emph{Say it:} \emph{kashu}, once more
  \item \emph{Say it:} say \emph{kashu}
  \item \emph{Recall:} say the Japanese for a farmhouse, then the Japanese for a carpenter, then say \emph{kashu} again
\end{itemize}

\begin{tcolorbox}[breakable,colback=teal!4,colframe=teal!35!black,title={Before you move on}]
What does \ja{かしゅ} mean? (\textquotedblleft{}A singer\textquotedblright{}.) Say it once more.
\end{tcolorbox}
\section[\texorpdfstring{\ja{しごと} (shigoto)}{しごと (shigoto)}]{\ja{しごと} (shigoto) — work, a job}

\label{lesson:JA-C112-shigoto}



\begin{quote}
Before the new one: say the Japanese for a carpenter, then the Japanese for a singer.
\end{quote}

\subsection*{You'll want to know: \ja{しごと}}
\textbf{\ja{しごと}} — \emph{shigoto} — \textquotedblleft{}work, a job\textquotedblright{}.

\begin{grammarlens}[title={What you've built}]
One more word for this chapter.
\end{grammarlens}

\subsection*{Your turn}
\begin{itemize}
\raggedright
  \item \emph{Say it:} \emph{shigoto}
  \item \emph{Say it:} \emph{shigoto}, once more
  \item \emph{Say it:} say \emph{shigoto}
  \item \emph{Recall:} say the Japanese for a carpenter, then the Japanese for a singer, then say \emph{shigoto} again
\end{itemize}

\begin{tcolorbox}[breakable,colback=teal!4,colframe=teal!35!black,title={Before you move on}]
What does \ja{しごと} mean? (\textquotedblleft{}Work, a job\textquotedblright{}.) Say it once more.
\end{tcolorbox}
\section[\texorpdfstring{\ja{やすみ} (yasumi)}{やすみ (yasumi)}]{\ja{やすみ} (yasumi) — a holiday, a break}

\label{lesson:JA-C112-yasumi}



\begin{quote}
Before the new one: say the Japanese for a singer, then the Japanese for work.
\end{quote}

\subsection*{You'll want to know: \ja{やすみ}}
\textbf{\ja{やすみ}} — \emph{yasumi} — \textquotedblleft{}a holiday, a break\textquotedblright{}.

\begin{grammarlens}[title={What you've built}]
One more word for this chapter.
\end{grammarlens}

\subsection*{Your turn}
\begin{itemize}
\raggedright
  \item \emph{Say it:} \emph{yasumi}
  \item \emph{Say it:} \emph{yasumi}, once more
  \item \emph{Say it:} say \emph{yasumi}
  \item \emph{Recall:} say the Japanese for a singer, then the Japanese for work, then say \emph{yasumi} again
\end{itemize}

\begin{tcolorbox}[breakable,colback=teal!4,colframe=teal!35!black,title={Before you move on}]
What does \ja{やすみ} mean? (\textquotedblleft{}A holiday, a break\textquotedblright{}.) Say it once more.
\end{tcolorbox}
